% =====================================================================
% CSE 341: Concepts of Programming Languages
% Project D4 (Part 2): AI Usage Journal
% Student: 220104004926 Yusuf Hamza Çelebi
% Compile with: pdflatex 220104004926_D4_P2.tex
% =====================================================================
\documentclass[11pt,a4paper]{article}

\usepackage[utf8]{inputenc}
\usepackage[T1]{fontenc}
\usepackage{lmodern}
\usepackage[english]{babel}
\usepackage{microtype}
\usepackage[a4paper,margin=2.25cm,top=2.0cm,bottom=2.15cm]{geometry}

\usepackage[table]{xcolor}
\definecolor{ansblue}{HTML}{1E3A8A}
\definecolor{anscaption}{HTML}{6B7280}
\definecolor{anscodebg}{HTML}{F9FAFB}
\usepackage[colorlinks=true, linkcolor=ansblue, urlcolor=ansblue,
            citecolor=ansblue, pdfborder={0 0 0}]{hyperref}

\usepackage[explicit]{titlesec}
\titleformat{\section}
  {\large\bfseries\sffamily\color{black}}
  {\thesection}{0.65em}{#1}
\titlespacing*{\section}{0pt}{13pt}{5pt}
\setlength{\parskip}{4.5pt plus 1pt minus 1pt}
\setlength{\parindent}{0pt}
\linespread{1.06}
\emergencystretch=2em
\setcounter{secnumdepth}{0}

\usepackage{amssymb}
\usepackage{enumitem}
\setlist[itemize]{leftmargin=1.1em,itemsep=1pt,topsep=2pt,parsep=0pt}

\usepackage{fancyhdr}
\pagestyle{fancy}
\fancyhf{}
\renewcommand{\headrulewidth}{0pt}
\fancyfoot[L]{\footnotesize\color{anscaption}%
  CSE 341: AI Usage Journal (D4, Part 2) \(\cdot\)
  Student 220104004926 Yusuf Hamza \c{C}elebi}
\fancyfoot[R]{\footnotesize\color{anscaption}Page \thepage}

\hypersetup{
  pdftitle  = {AI Usage Journal (D4, Part 2)},
  pdfauthor = {220104004926 Yusuf Hamza Celebi},
  pdfsubject= {CSE 341 Project Part 2},
  pdfkeywords={AI, usage journal, ANS, DSL, type system, interpreter}
}

\def\UrlBreaks{\do\/\do-\do_\do.\do~}
\newcommand{\code}[1]{\texttt{#1}}
\newcommand{\field}[2]{\textbf{#1.} #2\par}
\newcommand{\prompt}[1]{%
  \textbf{Prompt.}
  \begin{quote}\small #1\end{quote}
}
\newcommand{\sessionlog}[1]{\textbf{Session log.} \url{#1}\par}

\begin{document}

\begin{center}
  {\footnotesize\color{anscaption}CSE 341: Concepts of Programming Languages\par}
  \vspace{3pt}
  {\huge\bfseries\sffamily AI Usage Journal\par}
  \vspace{4pt}
  {\large Deliverable 4 (D4), Part 2\par}
  \vspace{6pt}
  {\footnotesize\color{anscaption}
    Student: 220104004926 Yusuf Hamza \c{C}elebi \(\cdot\)
    Spring 2026 \(\cdot\)
    Gebze Technical University\par}
\end{center}
\vspace{8pt}

% =====================================================================
\section{Entry 1 of 10}
\field{Date}{2026-05-08}
\field{Phase}{Writing / Journal setup}
\field{AI tool}{Claude Sonnet 4.6, claude.ai (desktop)}
\field{Goal}{Find where Claude Code stores session logs on macOS so all interactions can be retrieved for the journal.}
\prompt{``Where can I find the chat history of claude code chats in macos? Where are they located?''}
\field{Response key part}{\texttt{\textasciitilde/.claude/projects/} --- each project subfolder is named after the project's absolute path with slashes replaced by hyphens. Session files are \texttt{.jsonl}, one per conversation, named with a timestamp and UUID.}
\field{Accepted}{The path is correct. Confirmed by running \code{ls \textasciitilde/.claude/projects/} in Terminal and finding the \texttt{-Users-yhcelebi-Desktop-ans} subfolder with all session files.}
\field{Rejected / modified}{Nothing to reject. The answer was direct and verifiable.}
\field{Errors caught}{None. Independent verification with \code{find \textasciitilde/.claude/projects/ -name "*.jsonl" | sort} confirmed the complete file list.}
\field{Reflection}{Knowing the log location makes it straightforward to retrieve every session for the journal without relying on manual exports or memory. Running \code{find} gives a complete chronological list of every interaction across the whole project. This entry was drafted from Claude Desktop; a shareable link will be provided if the tutor requests direct access to the conversation.}
\sessionlog{https://claude.ai/share/e31ce86d-3082-40e4-896a-da16422320dc}

% =====================================================================
\section{Entry 2 of 10}
\field{Date}{2026-05-14}
\field{Phase}{Documentation / Writing}
\field{AI tool}{Claude Code, Claude Sonnet 4.6 CLI}
\field{Goal}{Produce a comprehensive \texttt{README.md} for the ANS repository that covers the language, the build system, and the project structure.}
\prompt{``Can you write a comprehensive README.md that explains and adds up all about ANS lang''}
\field{Response key part}{Before writing, Claude explored the full repository tree: D1 specification, D2 source files, D3 test programs, and sample \texttt{.ans} files. The README it produced covered the project overview, build and run instructions (\code{make}, \code{./ans -{}-run}, \code{./ans -{}-check}), a language reference with syntax examples, and an architecture section describing the lexer\,/\,parser\,/\,type-checker\,/\,interpreter pipeline.}
\field{Accepted}{The structure and build instructions were accurate against the actual Makefile and D2 layout. The language-reference examples matched the current grammar.}
\field{Rejected / modified}{A paragraph describing \code{-{}-dump-ast} as the primary usage mode was rewritten. For Part 2, \code{-{}-run} is the representative command, not \code{-{}-dump-ast}, which was the Part 1 entry point.}
\field{Errors caught}{The README initially listed \texttt{mathpazo} and \texttt{zi4} as document compilation dependencies. Those fonts only work under \texttt{tectonic}, not \texttt{pdflatex}. The D1\,P2 and D3\,P2 documents were compiled with \texttt{pdflatex} using \texttt{lmodern}, so the README dependency list was wrong and had to be corrected.}
\field{Reflection}{Writing documentation before finishing the interpreter forced an inventory of which features were specified but not yet implemented. Two gaps appeared: the \code{-{}-run} flag did not yet exist at the time of writing, and the README's promise that the checker rejects unbalanced reactions was ahead of the actual implementation state.}
\sessionlog{/Users/yhcelebi/Desktop/ans/Archive/YHC/claude_code/5fb50579-254d-47c5-89b4-6c910afebcc5.jsonl}

% =====================================================================
\section{Entry 3 of 10}
\field{Date}{2026-05-19}
\field{Phase}{Version control / Collaboration}
\field{AI tool}{Claude Code, Claude Sonnet 4.6 CLI}
\field{Goal}{Resolve a \code{git pull} failure on the \code{demo/enes} branch caused by divergent history.}
\prompt{``I cannot git pull demo/enes branch can you solve it''}
\field{Response key part}{Claude ran \code{git branch -vv} and \code{git log -{}-oneline} on both sides. It found that local and remote \texttt{demo/enes} had diverged: four commits existed on the remote that were not local, and four local commits were not on the remote. Claude merged the branches, then resolved two conflicts: (1) \texttt{README.md} --- both sides had modified it; (2) the old \texttt{D2 P1 zip} --- deleted locally but still present remotely. Claude deferred the README conflict to human judgment and explained the implication of each option before waiting.}
\field{Accepted}{The diagnostic approach (comparing remote and local commit counts before attempting any merge) and the decision to keep Enes's remote README version, which was more recent.}
\field{Rejected / modified}{Claude's first proposal was \code{git rebase}. Rebasing public branch history would rewrite commits that the partner had already built on, which breaks collaborative workflows. A plain merge was used instead.}
\field{Errors caught}{The initial response said ``no tracking remote branch set.'' That was incorrect: \code{git branch -vv} already showed the tracking set. Claude corrected this diagnosis on the next message after being shown the output directly.}
\field{Reflection}{AI was effective at reading git command output and forming a diagnosis quickly. However, conflict resolution still required human judgment about which version of \texttt{README.md} to keep. That decision depended on knowing the timeline of changes, which neither the AI nor a diff could determine on its own.}
\sessionlog{/Users/yhcelebi/Desktop/ans/Archive/YHC/claude_code/37e46d72-c73b-4290-9802-387b2f4e9791.jsonl}

% =====================================================================
\section{Entry 4 of 10}
\field{Date}{2026-05-20}
\field{Phase}{Design / Writing}
\field{AI tool}{Claude Code, Claude Sonnet 4.6 CLI}
\field{Goal}{Write the complete \S4.5 Type System section of D1 Part 2 following Sebesta Chapter\,6.}
\prompt{``Can you prepare a type system md file in D1/4.5 according to ANS and Sebesta''}
\field{Response key part}{Claude produced a full \S4.5 document covering: primitive types (Number, String, Character, Boolean), domain types (Quantity, Formula, List\texttt{<T>}, Element, Compound, Reaction, ReactionSide), strong typing, no implicit coercion, and type equivalence. It organised the section around five Sebesta Ch.\,6 subsections (\S6.2, \S6.5--6.7, \S6.14, \S6.15) and included a table for each structured-type decision.}
\field{Accepted}{The five-subsection organisation and Sebesta citation mapping. The primitive type value-set descriptions (IEEE-754 double for Number, Unicode with escapes for String) were accurate.}
\field{Rejected / modified}{The initial treatment of Boolean was incomplete --- it described Boolean as ``produced by comparisons; no literal'' but did not clarify that Boolean is not a valid \texttt{<type\_name>} parameter annotation. That distinction matters for the type checker.}
\field{Errors caught}{Claude placed Quantity inside the ``primitive types'' subsection in the first draft. Quantity is a domain type (a number--unit pair), not a primitive. Sebesta \S6.2 defines primitives as types provided by the language that are not defined in terms of other types. Quantity is built from Number and a unit tag, so it is a structured domain type.}
\field{Reflection}{AI drafts specification structure quickly but domain-specific nuance requires the author to review each claim against both the language spec and the textbook definition simultaneously. At least three review passes were needed to bring \S4.5 into full alignment with the actual ANS type rules.}
\sessionlog{/Users/yhcelebi/Desktop/ans/Archive/YHC/claude_code/d187f89a-c178-4fae-bee8-2f5da398a695.jsonl}

% =====================================================================
\section{Entry 5 of 10, Experiment E2}
\field{Date}{2026-05-20}
\field{Phase}{Design / Type System}
\field{AI tool}{Claude Code, Claude Sonnet 4.6 CLI}
\field{Goal}{Conduct Experiment E2: ask the AI to explain the difference between name equivalence and structural equivalence as defined in Sebesta \S6.15, apply the distinction to ANS record and list types, and provide a concrete example.}
\prompt{``The doc claims ANS lists are heterogeneous, so list[0] has no static type. But then calling ANS strongly typed and statically typed is a contradiction. Lists need to be either List\textless T\textgreater\ homogeneous or we drop the static typing claim. Explain name equivalence vs structural equivalence (Sebesta \S6.15) and how it applies here.''}
\field{Response key part}{Claude explained both rules from Sebesta \S6.15. \textbf{Name equivalence:} two types are identical only when they share the same declared type name --- \texttt{compound} matches \texttt{compound}; two compound values with identical formulas are still both of type \texttt{compound}, not of two distinct named types. \textbf{Structural equivalence:} two types are equal when their structures match recursively --- for \texttt{List<T>}, \texttt{List<Number>} and \texttt{List<Number>} are the same type because \texttt{Number} $\equiv$ \texttt{Number}. Claude proposed three rules for ANS: name\,/\,kind equivalence for Element, Compound, Reaction; unit-tag equality for Quantity; structural equivalence for \texttt{List<T>}.}
\field{Accepted}{The three-rule framework. It correctly maps to the Sebesta \S6.15 classification and resolves the heterogeneous\,/\,homogeneous question in favour of \texttt{List<T>} with structural equivalence.}
\field{Rejected / modified}{Claude initially suggested applying structural equivalence to Compound as well --- two compounds with the same formula would then be ``the same type'' at the checker level. That would break typed parameter checking: a function declared \texttt{function f(compound c)} would accept any compound value, which is the intended behaviour, but it would also mean the type of \texttt{water} and the type of \texttt{methane} are identical just because both are \texttt{compound}. The final rule is kind-based: all compound values share type \texttt{compound}, and typed parameters accept any compound value. Structural equivalence of formulas is not type equivalence.}
\field{Errors caught}{Claude's first explanation conflated structural equivalence of records (Compound, Element) with structural equivalence of parameterised collections (List\texttt{<T>}). In Sebesta's model these are distinct cases. Compound is a record-like type, so name equivalence applies; \texttt{List<T>} is parameterised, so structural equivalence of the type constructor is appropriate. I caught this by re-reading Sebesta \S6.15 alongside the checker.c type-string representation (\texttt{"compound"}, \texttt{"list<number>"}).}
\field{Reflection}{E2 was the most instructive AI exchange in Part 2. Forcing a precise Sebesta-sourced answer on type equivalence revealed an implicit design decision that the spec had left ambiguous. The AI gave the correct answer for List\texttt{<T>} but the wrong answer for Compound, and catching that difference required tracing the implications into the actual type-checker implementation.}
\sessionlog{/Users/yhcelebi/Desktop/ans/Archive/YHC/claude_code/d187f89a-c178-4fae-bee8-2f5da398a695.jsonl}

% =====================================================================
\section{Entry 6 of 10}
\field{Date}{2026-05-20}
\field{Phase}{Design / Type System}
\field{AI tool}{Claude Code, Claude Sonnet 4.6 CLI}
\field{Goal}{Determine whether a bare \code{[]} literal is valid in ANS and what the specification should say about empty list literals.}
\prompt{``The List\textless T\textgreater\ decision exposes a problem. If `let x = []`, what is T? ANS has no type annotation on `let`, so there is no context to fix T. What should the spec say?''}
\field{Response key part}{Claude identified the problem correctly: ANS's \texttt{let} binding infers the type purely from the right-hand side, and there is no \texttt{List<T>} annotation position in the grammar. Therefore \texttt{T} is unconstrained for a bare \texttt{[]}, and the language has no syntactic context that can fix the type parameter. Claude concluded that \texttt{[]} should be rejected as a semantic error and that every list literal must have at least one element so that \texttt{T} can be inferred from the element type.}
\field{Accepted}{The semantic rejection rule: \texttt{[]} is a semantic error, not a parse error. The spec now states this explicitly in the List\texttt{<T>} design-decisions table under ``Initialisation.'' This preserves the fully static element type without requiring a separate type-annotation syntax.}
\field{Rejected / modified}{The suggestion to allow \texttt{[]} with a special ``empty list'' type and defer element-type checking to the first indexing operation was rejected. That would require a form of deferred type inference that ANS does not otherwise support, and it would weaken the claim that every type is determined at the point of declaration.}
\field{Errors caught}{Claude initially proposed ``infer \texttt{T} from the assignment target type'' as a possible fix. ANS \texttt{let} has no declared type on the left side, so there is no target type. This was a generic type-inference suggestion that did not account for ANS's actual binding rules.}
\field{Reflection}{The empty-list case is small but matters for specification completeness. The AI surfaced it quickly by following the \texttt{List<T>} decision to its logical conclusion, but the correct fix required knowing what ANS deliberately does \emph{not} have (type annotations on variables) as much as what it does have.}
\sessionlog{/Users/yhcelebi/Desktop/ans/Archive/YHC/claude_code/d187f89a-c178-4fae-bee8-2f5da398a695.jsonl}

% =====================================================================
\section{Entry 7 of 10}
\field{Date}{2026-05-21}
\field{Phase}{Writing / Design Rationale}
\field{AI tool}{Claude Code, Claude Sonnet 4.6 CLI}
\field{Goal}{Write \S4.8 Design Rationale with five decisions from \S4.5--4.7, each citing a Sebesta Ch.\,5--7 section and following a choice\,$\to$\,reason\,$\to$\,trade-off structure.}
\prompt{``Write section 4.8 Design Rationale. Five decisions from 4.5--4.7. Each must follow the Sebesta Ch.\,5--7 framework: state the decision, cite the Sebesta section, give the reason in terms of readability / writability / reliability, state the trade-off.''}
\field{Response key part}{Claude produced five paragraphs: (1) static strong typing with no coercion (\S6.14, \S7.4); (2) Quantities as a domain type (\S6.2); (3) static lexical scoping (\S5.5); (4) left-to-right operand evaluation (\S7.2); (5) assignment as a statement, not an expression (\S7.7). Each paragraph closed with a cost or trade-off sentence.}
\field{Accepted}{The choice\,$\to$\,reason\,$\to$\,trade-off three-part structure and all five Sebesta section citations. Paragraphs 1, 3, and 5 were the strongest: each cited a concrete design criterion from Sebesta (reliability, readability) and named a specific trade-off.}
\field{Rejected / modified}{Paragraph 2 described Quantities as ``similar to physical units in other DSLs'' without citing a specific Sebesta concept. It was rewritten to reference \S6.2 with a concrete justification: a value such as \code{890.3 kJ/mol} is not interchangeable with a plain number, so preserving the unit tag inside the value makes the type checker aware of the distinction. Also, a missing edge case was added: same-unit addition (\code{5 kJ/mol + 890.3 kJ/mol}) is not defined in this subset even though the units match, and the spec should say so explicitly.}
\field{Errors caught}{The trade-off sentence in each paragraph was too vague in the first draft. For example, paragraph 4 said ``the trade-off is that evaluation order is fixed.'' That is not a trade-off --- it is the property. The real trade-off is that the interpreter cannot reorder operand evaluation for optimisation. Each trade-off sentence was rewritten by hand.}
\field{Reflection}{The three-part structure is harder to apply than it looks. The ``trade-off'' part forces the author to name something the design decision costs, and AI tends to produce statements that describe the property rather than the cost. The rationale paragraphs became genuinely useful only after each trade-off sentence was rewritten.}
\sessionlog{/Users/yhcelebi/Desktop/ans/Archive/YHC/claude_code/1221c730-cc2e-4725-a94d-14777f23c8ef.jsonl}

% =====================================================================
\section{Entry 8 of 10}
\field{Date}{2026-05-22}
\field{Phase}{Specification / Semantics}
\field{AI tool}{Claude Code, Claude Sonnet 4.6 CLI}
\field{Goal}{Fix a consistency error between D1 \S4.4 semantics and the D3 example programs: the spec only allowed \texttt{for-in} over a Formula, but D3 programs iterate over a ReactionSide.}
\prompt{``The D3 example program uses \texttt{for (let species in reactants)} where \texttt{reactants} is a ReactionSide. But \S4.4 currently only says the for-in iterable must be a Formula. Confirm this is a spec bug and propose the correct fix.''}
\field{Response key part}{Claude confirmed the inconsistency by quoting the \S4.4 rule and the D3 program side by side. It identified two possible fixes: (a) restrict D3 programs to iterate only over formulas; (b) extend \S4.4 to include ReactionSide as a valid iterable, specifying that the loop variable is bound to each term's species in declaration order, with coefficients accessible via \texttt{reaction\_side[compound]} rather than through the loop variable. Claude recommended (b) because the pattern is chemically natural and already implemented in the interpreter.}
\field{Accepted}{Fix (b) --- the \S4.4 semantics were extended with a ReactionSide iteration rule. The rule states that each iteration binds the loop variable to the species value (a Compound when the species name matches a declared compound), and that iteration order follows the declaration order of terms.}
\field{Rejected / modified}{Claude's first suggestion included also allowing \texttt{List<T>} in the same \S4.4 extension. That was rejected as out of scope: D3 programs do not iterate over lists, and the change would require additional type-checking rules not yet covered.}
\field{Errors caught}{Claude initially proposed fixing the D3 programs to use formula iteration rather than extending the spec. That was the wrong direction: the D3 programs had already been run through the actual interpreter and produced correct output. The spec was the stale document, not the programs.}
\field{Reflection}{Maintaining consistency between a specification, an implementation, and test programs across two submission phases is the hardest editorial problem in a multi-deliverable project. AI found the inconsistency quickly by reading all three artefacts simultaneously, which would have taken much longer as a manual cross-check.}
\sessionlog{/Users/yhcelebi/Desktop/ans/Archive/YHC/claude_code/597d86e4-d5ec-4881-974c-d9c1ebd43d7f.jsonl}

% =====================================================================
\section{Entry 9 of 10}
\field{Date}{2026-05-22}
\field{Phase}{Implementation / Bug Fix}
\field{AI tool}{Claude Code, Claude Sonnet 4.6 CLI}
\field{Goal}{Confirm and fix a reported parser bug where \texttt{1 + 2 * 3;} fails to parse as an expression statement.}
\prompt{``A peer audit found: expression statement parsing is broken --- \texttt{1 + 2 * 3;} fails at the \texttt{*} token. Confirm and fix.''}
\field{Response key part}{Claude reproduced the failure and traced it to two sites in \texttt{parser.c}. In \texttt{parse\_top\_level}, the additive loop called \texttt{parse\_unary(p)} as its right operand instead of \texttt{parse\_multiplicative(p)}. For the input \texttt{1 + 2 * 3}, this consumed \texttt{2} and left \texttt{* 3} unconsumed, causing the subsequent \texttt{expect(;)} to fail. A second site in \texttt{parse\_statement} had the multiplicative and additive loops in the wrong order (additive ran before multiplicative). Both sites needed to be corrected.}
\field{Accepted}{Both fixes: correcting the loop order in \texttt{parse\_statement} and changing \texttt{parse\_unary(p)} to \texttt{parse\_multiplicative(p)} as the additive right operand in both \texttt{parse\_top\_level} and \texttt{parse\_statement}. The 35-test suite passed cleanly after both changes.}
\field{Rejected / modified}{Nothing. The analysis was correct and the fix was minimal.}
\field{Errors caught}{Claude's first explanation named only \texttt{parse\_top\_level} as the problem. When prompted, it found the same issue in \texttt{parse\_statement}. An incomplete single-site fix would have left the bug present for any expression statement inside a block.}
\field{Reflection}{The existing test suite did not exercise \texttt{1 + 2 * 3;} as a standalone expression statement. Tests covered expressions inside conditions and variable declarations, but not the expression-statement parse path directly. A grammar-focused test suite can miss bugs that only surface in a specific statement position.}
\sessionlog{/Users/yhcelebi/Desktop/ans/Archive/YHC/claude_code/597d86e4-d5ec-4881-974c-d9c1ebd43d7f.jsonl}

% =====================================================================
\section{Entry 10 of 10, Experiment E3}
\field{Date}{2026-05-22}
\field{Phase}{Implementation / Type Checker}
\field{AI tool}{Claude Code, Claude Sonnet 4.6 CLI}
\field{Goal}{Conduct Experiment E3: evaluate the AI-generated type checker on three concrete test cases, identify the bug reported in a peer audit, and verify the fix.}
\prompt{``Audit finding: \texttt{check\_function\_decl} calls \texttt{tenv\_new(NULL)}, creating an isolated type environment with no parent. Functions cannot see top-level \texttt{let} variables at check time, even though the interpreter finds them via \texttt{g\_global\_env}. Confirm with test cases and fix.''}
\field{Response key part}{Claude read \texttt{checker.c} and \texttt{interpreter.c} in parallel and confirmed the asymmetry. The checker created function scopes with \texttt{tenv\_new(NULL)} (no parent), while the interpreter created \texttt{Env} frames with \texttt{parent = g\_global\_env}. Any function referencing a top-level \texttt{let} variable would receive an ``undeclared name'' type error at check time but run without error at runtime. The fix was to pass \texttt{global\_tenv} to \texttt{check\_function\_decl} and use \texttt{tenv\_new(global\_tenv)} as the function frame.}

\textbf{Test case 1} (top-level \texttt{let} referenced inside function):
\begin{quote}\small
\code{let n = 20.0; function f() \{ return n; \};}
\end{quote}
Expected: no type error (\texttt{n} visible by static scoping). Before fix: checker reports \texttt{n} undeclared. After fix: no error. \checkmark

\textbf{Test case 2} (top-level compound referenced inside function body):
\begin{quote}\small
\code{compound water \{ formula = H2O, name = "Water" \};}\\
\code{function f(compound c) \{ let x = water.name; return x; \};}
\end{quote}
Expected: no error. Before fix: checker reports \texttt{water} undeclared inside \texttt{f}. After fix: no error. \checkmark

\textbf{Test case 3} (local shadowing must still work after fix):
\begin{quote}\small
\code{let n = 20.0; function f() \{ let n = 99.0; return n; \};}
\end{quote}
Expected: no error; local \texttt{n} shadows outer \texttt{n}. Before and after fix: no error. \checkmark

\field{Accepted}{The fix --- \texttt{tenv\_new(global\_tenv)}. This mirrors the interpreter's scoping model exactly: functions can see all global declarations but not each other's local frames.}
\field{Rejected / modified}{Claude's broader suggestion to ``pass the caller environment as the parent'' was rejected. That would implement dynamic scoping (each call site's locals visible to the callee), not the static scoping defined in \S4.6 and implemented in the interpreter. The fix must use the fixed global environment, not the current call environment.}
\field{Errors caught}{Claude initially suggested that \texttt{tenv\_new(NULL)} might be intentional, to prevent functions from capturing outer variables as closures. That reading was wrong. Closures require capturing the \emph{definition-site} environment; static scoping requires using the \emph{global} environment as parent. These are different things: ANS has static scoping but not closures.}
\field{Reflection}{E3 produced the most practically significant finding. A type checker that disagrees with the interpreter on scoping is worse than having no type checker: it rejects valid programs at check time while the interpreter runs them correctly. The bug was subtle because the tests that existed all called functions from the top level with no top-level \texttt{let} references inside function bodies. The AI found the mismatch by comparing the two modules simultaneously, which is the kind of cross-file analysis that is easy to miss in incremental development.}
\sessionlog{/Users/yhcelebi/Desktop/ans/Archive/YHC/claude_code/597d86e4-d5ec-4881-974c-d9c1ebd43d7f.jsonl}

\end{document}
